\documentclass[border=10pt,crop=true]{standalone}
\usepackage{circuitikz}
% Shared style for 电子学一 Chapter 2 op-amp circuit figures (CircuiTikZ).
\usetikzlibrary{calc}

\ctikzset{
  bipoles/length=0.95cm,
  bipoles/thickness=1,
  resistors/scale=1,
  capacitors/scale=1,
  label distance=2.5mm,
}

\tikzset{
  opfig/.style={font=\small},
  opfigThin/.style={font=\scriptsize},
}

% Geometry (cm) — keep identical across all op-amp schematics
\def\OpInLen{1.55}    % label → input terminal
\def\OpInGap{0.08}    % terminal → wire to op amp +
\def\OpAmpWire{1.00}  % terminal side → op amp + anchor
\def\OpOutLen{1.05}   % op amp out → output terminal
\def\OpJuncL{0.50}    % v_- → feedback junction (left)
\def\OpFbDown{0.55}   % v_- straight down before R_1
\def\OpRvert{1.20}    % R_1 vertical span to ground
\def\OpInSep{0.55}    % dual-input spacing (v_1 / v_2)
\def\OpFbStub{0.40}   % follower: stub right of output before dropping
\def\OpFbClear{0.22}  % follower: below op-amp south for horizontal feedback


\begin{document}
\begin{circuitikz}[american, scale=1.0, transform shape]
  \draw (-\OpInLen, 0) node[left, opfig]{$v_+$} to[short, o-] ++(\OpInGap, 0)
    -- ++(\OpAmpWire, 0)
    node[op amp, noinv input up, anchor=+](OA){}
    (OA.out) to[short, -o] ++(\OpOutLen, 0) node[right, opfig]{$v_o$}
    (OA.-) -- ++(-\OpJuncL, 0) node[ground, scale=0.85]{};

  \node[opfigThin, anchor=north] at ($(OA.-)+(-0.26, -0.04)$) {$v_-$};

  \draw (OA.up) -- ++(0, 0.48);
  \node[opfig, anchor=south] at ($(OA.up)+(0, 0.52)$) {$V_{cc}$};
  \node[opfigThin, anchor=south] at ($(OA.up)+(0, 0.30)$) {$+10{\sim}15\,\mathrm{V}$};

  \draw (OA.down) -- ++(0, -0.48);
  \node[opfig, anchor=north] at ($(OA.down)+(0, -0.52)$) {$V_{ee}$};
  \node[opfigThin, anchor=north] at ($(OA.down)+(0, -0.30)$) {$-10{\sim}{-15}\,\mathrm{V}$};
\end{circuitikz}
\end{document}
